\documentclass[../main]{subfiles}
\begin{document}
\section{Repaso}
\begin{center}
\includegraphics[width=0.2\textwidth]{images/dibujotecnico.png}
\end{center}
\subsection{Revisión:Escalas, líneas figuras y caligrafía.}
\begin{enumerate}
\item Dibujar las siguientes piezas en escala 1:1, respetando los trazos en los tipos de línea, las medidas están en milímetros. Completar en  la siguiente página actividad de caligrafía.
\end{enumerate}
\begin{center}
\includegraphics[width=0.95\textwidth]{images/06/dibujo1.png}
\end{center}

\newpage
\begin{tikzpicture}[remember picture, overlay, shift={(-3.6,-27.8)}]
\tikzstyle{every node}=[font=\Huge]
\draw [step=0.1,gray,line width=0.1mm] (2.5,4.5) grid (20,28.7) ;
\draw [step=1cm,black,line width=0.1mm] (2.5,4.5) grid (20,28.7) ;

\node at (4,28.2) [fontscale=4]{a a a a};
\node at (3,27.3) [fontscale=4]{b};
\node at (3,26.2) [fontscale=4]{c};
\node at (3,25.3) [fontscale=4]{d};
\node at (3,24.2) [fontscale=4]{e};
\node at (3,23.3) [fontscale=4]{f};
\node at (3,22.2) [fontscale=4]{x};
\node at (3,21.3) [fontscale=4]{h};
\node at (3,20.3) [fontscale=4]{i};
\node at (3,19.3) [fontscale=4]{j};
\node at (3,18.3) [fontscale=4]{k};
\node at (3,17.3) [fontscale=4]{l};
\node at (3,16.2) [fontscale=4]{m};
\node at (3,15.2) [fontscale=4]{n};
\node at (3,14.2) [fontscale=4]{o};
\node at (3,13.1) [fontscale=4]{p};
\node at (3,12.1) [fontscale=4]{q};
\node at (3,11.2) [fontscale=4]{z};
\node at (3,10.3) [fontscale=4]{t};
\node at (3,9.2) [fontscale=4]{u};
\node at (3,8.2) [fontscale=4]{v};
\node at (3,7.2) [fontscale=4]{w};
\node at (3,6.2) [fontscale=4]{x};
\node at (3,5.3) [fontscale=4]{\%};
\end{tikzpicture}


\end{document}